% Part: history
% Chapter: biographies 
% Section: alonzo-church

\documentclass[../../../include/open-logic-section]{subfiles}

\begin{document}

\olfileid{his}{bio}{chu}

\olsection{Alonzo Church}

\olphoto{church-alonzo}{Alonzo Church}

Alonzo Church（阿隆佐·丘奇）于1903年6月14日出生在华盛顿特区。幼年时，一次气枪事故导致他一只眼睛失明。他于1920年在康涅狄格州完成预备学校教育，并于同年进入普林斯顿大学开始大学学业，于1927年完成博士学业。在国外度过几年后，丘奇回到了普林斯顿。丘奇以极其礼貌和严谨而闻名。他的板书一丝不苟，并且会用 Duco cement（一种透明胶水）仔细涂抹重要文件以作保存。在学术研究之外，他喜欢阅读科幻杂志，如果发现文中有任何不严谨之处，他会毫不避讳地给编辑写信。

丘奇的学术成就卓著。他与学生 Stephen Kleene（克林）和 Barkley Rosser（罗瑟）一道，独立于 Alan Turing（艾伦·图灵）发明图灵机的工作，提出了一个关于能行可计算性的理论——λ演算。这两种可计算性的定义是等价的，并引出了现在所谓的\emph{丘奇–图灵论题}，即一个自然数函数是能行可计算的，当且仅当它是图灵机（或λ演算）可计算的。他还证明了现在所谓的\emph{丘奇定理}：一阶公式有效性的判定问题是不可解的。

丘奇步入高龄后仍继续工作。1967年他离开普林斯顿前往加州大学洛杉矶分校（UCLA），在那里担任教授直至1990年退休。丘奇于1995年8月1日逝世，享年92岁。

\begin{reading}
关于丘奇的简要生平，参见 \citet{EndertonND}。丘奇关于λ演算和判定问题（丘奇论题）的原始著述是 \citet{Church1936,Church1936a}。\citet{Aspray1984} 记录了对丘奇的一次访谈，内容关于1930年代的普林斯顿数学界。丘奇从1936年至1979年为《符号逻辑杂志》撰写了一系列书评。它们都存档在 John MacFarlane 的网站上 \citep{MacFarlane2015}。
\end{reading}

\end{document}